\documentclass[12pt]{article}
\usepackage{amsmath,amsthm,amssymb,amsfonts,setspace}
\usepackage{fancyhdr}
\usepackage{hyperref}

\textwidth 7.0 truein
\oddsidemargin -0.25in
\evensidemargin -0.5 truein
\topmargin -0.85in
\textheight 9.5in

%%%% Custom commands %%%%

% Box/square the final answer
\newcommand{\ans}[1]{\boxed{#1}}

% Shared indentation lengths for subproblems and solutions
\newlength{\subproblemindent}
\setlength{\subproblemindent}{1.5em}
\newlength{\subproblemlabelwidth}
\setlength{\subproblemlabelwidth}{2em}
\newlength{\solutionindent}
\setlength{\solutionindent}{3em}

% Make \circ usable in ordinary text as well as math mode
\let\latexcirc\circ
\renewcommand{\circ}{\ifmmode\latexcirc\else\ensuremath{\latexcirc}\fi}

% Bold a top-level problem: number and text separated by an em-dash
\newcommand{\problem}[2]{\vspace{1em}\noindent\textbf{#1} --- #2\par\vspace{0.25em}}

% Bold and indent a subproblem: letter and text separated by an em-dash
\newcommand{\subproblem}[2]{%
  \par\vspace{0.5em}%
  \noindent
  \hangindent=\dimexpr\subproblemindent+\subproblemlabelwidth+0.5em\relax
  \hangafter=1
  \hspace*{\subproblemindent}%
  \makebox[\subproblemlabelwidth][l]{\textbf{#1}}#2\par
}

% Indent all lines of the answer/solution block
\newcommand{\solution}[1]{%
  \par\vspace{0.0em}%
  \begin{list}{}{
    \setlength{\leftmargin}{\solutionindent}
    \setlength{\rightmargin}{0pt}
    \setlength{\listparindent}{0pt}
    \setlength{\itemindent}{0pt}
    \setlength{\parsep}{0pt}
  }%
    \item\relax
    \begingroup
      \setlength{\baselineskip}{1.5\baselineskip}
      \setlength{\lineskip}{1.5\lineskip}
      \setlength{\parskip}{0pt}
      \everymath{\displaystyle}
      #1
      \par
    \endgroup
  \end{list}%
  \vspace{0.25em}%
}

% Running footer with a link back to the source repo, shown on every page
\pagestyle{fancy}
\fancyhf{}
\renewcommand{\headrulewidth}{0pt}
\renewcommand{\footrulewidth}{0.4pt}
\fancyfoot[C]{\footnotesize\textit{Source: \url{https://github.com/chrismcarrico/homework}}}

%%%% In most cases you won't need to edit anything above this line %%%%

\begin{document}
\hfill Christian Carrico

\hfill SPCE 5085

\hfill Assignment \#3

\problem{1}{Calculate the system noise temperature of a 2 GHz receiver, given the following:  Tin=20 K, Trf=45, Tm = 450 K, Tif =800 K,  Gm=1 dB,  Grf = 20 dB, G if = 30 dB. 
See text section 4.3.2  (15 pts)   ; Where Tin = Temperature In, Trf = Temperature of RF elements, Tm = Temperature of mixer, Temperature of Intermediate frequency, and Gm = gain of mixer, Gif = gain of IF.    
Reminder: Use consistent units: don't mix dB and linear, convert dB units to linear before using equation 4.20.}
\solution{}
\problem{2}{Construct a graph showing the Doppler frequency (KHz) for a satellite in LEO (1000 km altitude) transmitting in L band microwave (1.5 GHz) from a ground station at observation angle of 10-170 degrees.  
A spreadsheet at 5 degree increments works well.  (15 pts)  Suggest that you read  article posted in module 2 on Doppler in LEO.}
\solution{}
\problem{3}{Construct a complete spread sheet downlink budget, for the Ku-band DBS-TV (Table 4.6a as is in textbook) 
Inserting all appropriate equations in the spread sheet to reproduce the values. (note typo in table, the path distance units should be 38,000 km (not meters).)   (40 pts)}
\solution{}
\subproblem{a}{Is the selection of QPSK at 1/2 rate an appropriate choice? explain}
\solution{}
\subproblem{b}{What would be the CNR and link margin be, if the Transponder was 170 W (vs. 160 W as in table) }
\solution{}
\problem{4}{Using 14 GHz: A Ku-band satellite uplink has a carrier frequency of 14.0 GHz and a carries a symbol stream at $R_{s} = 20 Msps$. The transmitter and receiver have SRRC filters with $\alpha = .20$}
\subproblem{a}{What is bandwidth occupied by the RF signal?}
\solution{}
\subproblem{b}{What is the frequency range of the transmitted RF signal?}
\solution{}
\problem{5}{Using Figure 5.13 on p.229 of my textbook, estimate graphically:}
\subproblem{a}{the minimum CNR (dB) modulated in BPSK to achieve BER of 10-6 or better.}
\solution{}
\subproblem{what is the expected BER for a QPSK modulated signal when CNR is about 15 dB}
\solution{}
\end{document}
